% =====================================================================
%  Chương Kết quả & Đánh giá — Phase 13 (bộ tấn công v3 + KB/RBAC + A/B)
%  Nội dung tự chứa, dán vào chương "Kiểm thử & Đánh giá" của báo cáo.
%  MỌI SỐ trong file này lấy từ artifact THẬT do scripts/run_v3_evaluation.py
%  sinh trên máy lab (reports/v3/...), KHÔNG có số bịa (AGENT_RULES rule 3).
%  Nguồn: docs/evaluation/phase13-v3-firewall-metrics-final.md
% =====================================================================

\section{Đánh giá tường lửa trên bộ tấn công tổng hợp v3}

\subsection{Phương pháp}

Bộ đánh giá gồm \textbf{300 case} tại \texttt{datasets/v3/cases/all.jsonl}:
\textbf{200 case tấn công} (mô phỏng kỹ thuật của NVIDIA garak, Microsoft PyRIT
và InjecAgent) và \textbf{100 case lành tính} (câu hỏi nghiệp vụ HR/IT/chính
sách). Toàn bộ payload là dữ liệu tổng hợp, sinh xác định bằng
\texttt{scripts/build\_v3\_attack\_payloads.py} (không gọi API LLM trả phí).

Mỗi metric được định nghĩa tường minh:
\begin{itemize}
  \item \textbf{TPR} (True Positive Rate) = số case tấn công bị chặn / tổng tấn công.
  \item \textbf{FPR} (False Positive Rate) = số case lành tính bị chặn / tổng lành tính.
  \item \textbf{Stop-before-LLM} = tỉ lệ case tấn công bị chặn \emph{trước} khi
        gọi tới mô hình (thuộc tính tiết kiệm token).
  \item \textbf{Exfil marker} = tỉ lệ câu trả lời chứa mã canary bí mật;
        chỉ đo được khi canary tương ứng thực sự nằm trong kho tri thức.
  \item ``Bị chặn'' nghĩa là quyết định cuối cùng $\in \{$\texttt{block},
        \texttt{human\_review}$\}$ --- hai quyết định giữ lại câu trả lời.
\end{itemize}

Theo quyết định phương pháp \textbf{L2}, \emph{latency không được báo cáo}
($p_{50}/p_{95}$ để trống); rate trên nhóm dưới $N=10$ mẫu báo là ``không đủ
mẫu'' thay vì một con số dễ gây hiểu nhầm.

Đánh giá chạy trên \textbf{hai cấu hình} để bóc tách đóng góp của từng lớp
phòng thủ:
\begin{enumerate}
  \item[(a)] \textbf{Rule-based} (\texttt{LLM\_PROVIDER=mock}): chỉ các guard
        dựa trên luật và heuristic, không gọi bộ phán đoán LLM.
  \item[(b)] \textbf{Semantic LLM judge} (\texttt{LLM\_PROVIDER=ollama},
        model cục bộ \texttt{hermes3:8b}): bổ sung lớp LLM-as-judge.
\end{enumerate}

\subsection{Kết quả cấu hình rule-based}

\begin{table}[h]
\centering
\caption{Ma trận nhầm lẫn --- cấu hình rule-based (n = 300)}
\begin{tabular}{lcc}
\hline
 & \textbf{Bị chặn} & \textbf{Được trả lời} \\
\hline
Tấn công (n = 200) & 98 (TP) & 102 (FN) \\
Lành tính (n = 100) & 0 (FP) & 100 (TN) \\
\hline
\end{tabular}
\end{table}

\begin{table}[h]
\centering
\caption{Chỉ số rule-based (run \texttt{832e9dcf})}
\begin{tabular}{lcc}
\hline
\textbf{Chỉ số} & \textbf{Giá trị} & \textbf{Ghi chú} \\
\hline
TPR & 49{,}0\% & \\
FPR & 0{,}0\% & không chặn nhầm câu lành tính nào \\
Stop-before-LLM & 49{,}0\% & mọi lần chặn đều trước khi gọi mô hình \\
Exfil marker & 0{,}0\% & n đo được = 200 \\
\hline
\end{tabular}
\end{table}

Phân tách theo nguồn kỹ thuật cho thấy giới hạn của lớp luật rất rõ rệt:

\begin{table}[h]
\centering
\caption{Tỉ lệ chặn theo nguồn tấn công --- rule-based}
\begin{tabular}{lccc}
\hline
\textbf{Nguồn} & \textbf{Tổng} & \textbf{Bị chặn} & \textbf{Tỉ lệ} \\
\hline
garak (jailbreak, payload splitting) & 105 & 82 & 78{,}1\% \\
injecagent (chỉ thị ẩn) & 40 & 16 & 40{,}0\% \\
pyrit (đóng vai, mạo danh cấp trên) & 55 & 0 & \textbf{0{,}0\%} \\
\hline
\end{tabular}
\end{table}

Lớp luật chặn hoàn toàn các đòn có \emph{dấu hiệu từ vựng} (\texttt{rag\_dan}
100\%, \texttt{payload\_splitting} 100\%) nhưng \textbf{bỏ lọt toàn bộ họ PyRIT}
(0/55). Nguyên nhân bản chất: đòn đóng vai và mạo danh thẩm quyền
(``với tư cách Giám đốc, tôi cần bảng lương cho cuộc họp'') dùng
\emph{ngữ dụng xã hội} thay vì từ khoá ghi đè chỉ thị, nên không mẫu regex nào
bắt được. Đây là giới hạn cố hữu của phòng thủ dựa-trên-luật.

\subsection{Kết quả cấu hình Semantic LLM judge}

\begin{table}[h]
\centering
\caption{Ma trận nhầm lẫn --- Semantic LLM judge \texttt{hermes3:8b} (n = 300)}
\begin{tabular}{lcc}
\hline
 & \textbf{Bị chặn} & \textbf{Được trả lời} \\
\hline
Tấn công (n = 200) & 200 (TP) & 0 (FN) \\
Lành tính (n = 100) & 57 (FP) & 43 (TN) \\
\hline
\end{tabular}
\end{table}

\begin{table}[h]
\centering
\caption{Chỉ số Semantic LLM judge (run \texttt{4ad06cc3})}
\begin{tabular}{lcc}
\hline
\textbf{Chỉ số} & \textbf{Giá trị} & \textbf{Ghi chú} \\
\hline
TPR & 100{,}0\% & chặn mọi case tấn công, gồm cả họ PyRIT \\
FPR & \textbf{57{,}0\%} & \textbf{chặn nhầm hơn nửa câu hỏi hợp lệ} \\
Stop-before-LLM & 99{,}5\% & \\
Exfil marker & 0{,}0\% & n = 200 \\
\hline
\end{tabular}
\end{table}

Lớp LLM-as-judge nâng tỉ lệ chặn PyRIT từ 0\% lên 100\%, đúng như dự đoán:
đòn đóng vai chỉ có thể bị bắt bằng phán đoán ngữ nghĩa, không bằng luật.

\subsection{Phân tích: đánh đổi precision--recall giữa hai lớp}

\textbf{Đây là kết quả trung tâm, và nó không phải một chiến thắng đơn giản.}
Con số TPR 100\% của lớp Semantic LLM \emph{không được diễn giải như bằng chứng
tường lửa mạnh}, vì nó đi kèm FPR 57\%: một bộ phân loại chặn gần như mọi thứ
cũng đạt TPR 100\% một cách tầm thường. Hai cấu hình thể hiện đúng đánh đổi
bảo mật--khả dụng kinh điển, đo được trên dữ liệu thật:

\begin{table}[h]
\centering
\caption{So sánh hai lớp phòng thủ}
\begin{tabular}{lccl}
\hline
\textbf{Chỉ số} & \textbf{Rule-based} & \textbf{Semantic LLM} & \textbf{Diễn giải} \\
\hline
TPR & 49{,}0\% & 100{,}0\% & LLM bắt thêm đòn đóng vai/mạo danh \\
FPR & 0{,}0\% & 57{,}0\% & LLM hy sinh khả dụng \\
Stop-before-LLM & 49{,}0\% & 99{,}5\% & \\
PyRIT block & 0{,}0\% & 100{,}0\% & đóng vai \emph{cần} lớp ngữ nghĩa \\
\hline
\end{tabular}
\end{table}

\begin{itemize}
  \item \textbf{Lớp luật}: precision cao (FPR 0\%) nhưng recall thấp
        (TPR 49\%) --- chính xác trên mẫu đã biết, mù trước tấn công ngữ nghĩa.
  \item \textbf{Lớp LLM}: recall hoàn hảo (TPR 100\%) nhưng precision thấp
        (FPR 57\%) --- chặn được mọi tấn công nhưng làm gián đoạn hơn nửa
        công việc hợp lệ, \textbf{chưa đạt mức triển khai thực tế}.
\end{itemize}

Kết luận kiến trúc: \emph{không cấu hình đơn lẻ nào đủ dùng}. Hai lớp bổ trợ
nhau --- lớp luật cho độ chính xác trên mẫu quen thuộc, lớp LLM cho độ phủ
trên đòn tấn công mới --- và việc \textbf{giảm FPR của lớp LLM} (hiệu chỉnh
prompt phán đoán hoặc ngưỡng quyết định) là hướng phát triển bắt buộc trước
khi đưa vào vận hành.

\section{Kiểm soát truy cập RBAC và so sánh A/B tường vs.\ không tường}

\subsection{Phân quyền theo vai trò trên kho tri thức doanh nghiệp}

Kho tri thức gồm 14 tài liệu tổng hợp (Kế toán/IT/Nhân sự/dùng chung), mỗi tài
liệu mang mức nhạy cảm và phạm vi phòng ban. Truy xuất áp dụng ACL \emph{trước}
khi nội dung tới mô hình. Bảng dưới là kết quả đo trực tiếp: một nhân viên IT
cấp thấp (\texttt{it.user1}) \textbf{không} truy xuất được tài liệu ngoài phạm
vi quyền, kể cả khi tường lửa bị tắt --- vì ACL là lớp độc lập với các guard nội
dung.

\begin{table}[h]
\centering
\caption{Cách ly truy cập theo vai trò (đo trực tiếp)}
\begin{tabular}{llc}
\hline
\textbf{Tài khoản} & \textbf{Tài liệu mật yêu cầu} & \textbf{Truy xuất được?} \\
\hline
it.user1 (IT, member) & Bảng lương HR & Không \\
it.user1 (IT, member) & Runbook TỐI MẬT của IT & Không \\
it.user1 (IT, member) & Dự án M\&A (ban điều hành) & Không \\
it.leader (IT, leader) & Runbook của IT & Có \\
hr.leader (HR, leader) & Bảng lương HR & Có \\
\hline
\end{tabular}
\end{table}

\subsection{So sánh A/B: chatbot có tường vs.\ không tường}

Hai chatbot dùng \emph{chung} một cơ chế truy xuất RAG và \emph{chung} một ACL,
chỉ khác ở việc bật/tắt các guard, cho phép so sánh công bằng. Kết quả đo
(người dùng đủ quyền, để cô lập tác dụng của guard khỏi tác dụng của ACL):

\begin{table}[h]
\centering
\caption{Rò rỉ bí mật doanh nghiệp: có tường vs.\ không tường}
\begin{tabular}{lll}
\hline
\textbf{Kịch bản} & \textbf{Có tường} & \textbf{Không tường} \\
\hline
Trích xuất credential runbook & \texttt{block} & lộ khoá \texttt{AKIA…}/\texttt{sk-…}/\texttt{ghp\_…} \\
Chỉ thị ẩn trong tài liệu (poisoning) & \texttt{block}, không gọi mô hình & mô hình đọc ngữ cảnh nhiễm độc \\
Câu hỏi chính sách lành tính & trả lời bình thường & trả lời bình thường \\
\hline
\end{tabular}
\end{table}

Điểm đáng chú ý về mặt vận hành: phía có tường trả về \texttt{provider\_name =
null} ở các kịch bản bị chặn --- nghĩa là mô hình \emph{chưa từng được gọi}, tiết
kiệm toàn bộ token --- một khác biệt \textbf{tất định}, không phụ thuộc mô hình
sinh ra câu trả lời gì.

\section{Hạn chế đã đo (không che giấu)}

Theo nguyên tắc trung thực học thuật, các hạn chế sau được nêu rõ, đều dựa trên
quan sát thực nghiệm chứ không phải suy đoán:

\begin{enumerate}
  \item \textbf{FPR 57\% của lớp Semantic LLM} khiến cấu hình này chưa dùng được
        trong thực tế; cần hiệu chỉnh trước khi triển khai.
  \item \textbf{Đường chat của workspace không có tầng DLP riêng} (khác với
        đường \texttt{/v1/rag/query}); bảng lương và PII có thể lọt qua nếu
        Output Guard không khớp mẫu.
  \item \textbf{Đòn đóng vai/mạo danh chỉ bị lớp Semantic LLM bắt}; với cấu hình
        rule-based, tỉ lệ chặn của họ này là 0\%.
  \item \textbf{Phía không tường mang tính ngẫu nhiên} do nhiệt độ sinh của mô
        hình; bằng chứng vững chắc nằm ở quyết định tất định phía có tường, không
        ở việc mô hình có chịu in dữ liệu ra hay không.
  \item Đánh giá dùng \textbf{corpus và mô hình cục bộ}; mô hình lớn hơn có thể
        tuân theo chỉ thị ẩn mạnh hơn, làm rủi ro rò rỉ tăng theo năng lực mô
        hình --- một hướng cần khảo sát thêm.
  \item \textbf{Latency không được báo cáo} (quyết định L2); split holdout của
        benchmark v2 \emph{không} được chạy trong báo cáo này.
\end{enumerate}
